% Fichier complété par tools/complete_missing_corrections.py.
% Source : CIN/CIN-01-Parametrage/1019_RobotPeinture/1019_RobotPeinture.tex
% Statut : rédigé (4 question(s)).
\begin{corrigebox}[Corrigé rédigé]
\CorrigeQuestion{1}
Une classe d'équivalence regroupe toutes les pièces sans mouvement relatif. On colorie donc d'une même couleur le bâti, d'une autre chaque sous-ensemble mobile et l'on sépare systématiquement les organes reliés par pivot ou glissière. Les éléments roulants, vis et écrous solidaires sont rattachés au solide qui les entraîne.
\CorrigeQuestion{2}
On obtient au minimum le bâti $S_0$, le sous-ensemble d'entrée $S_1$ et le sous-ensemble de sortie $S_2$; les éventuels renvois ou biellettes constituent chacun une classe supplémentaire. La liste exacte se lit en regroupant les repères de pièces qui ne sont séparés par aucune surface fonctionnelle de mouvement.
\CorrigeQuestion{3}
Le graphe comporte un sommet par classe d'équivalence et une arête par liaison. Les surfaces cylindriques coaxiales donnent un pivot, les guidages cylindriques avec translation un pivot glissant, et les guidages prismatiques une glissière. Les axes et centres sont ceux du dessin d'ensemble; les contacts de transmission sont ajoutés comme engrènement ou vis--écrou.
\CorrigeQuestion{4}
Le schéma cinématique minimal conserve le bâti, les axes des pivots, les directions des glissières et les longueurs utiles. Les dimensions de forme sans influence cinématique sont supprimées. Le nombre de mobilités du schéma doit être identique à celui du mécanisme réel et le graphe des liaisons doit être retrouvé exactement.
\end{corrigebox}
